\section{Respuesta a la RQ2}\label{sec:esquema:respuesta}\label{sec:evolucion:respuesta}

La RQ2 pregunta si ACM2YAML es suficientemente expresivo para representar la capacidad contractual de las APIs de correo, y qué extensiones requiere.
Con las ocho APIs del catálogo, la respuesta a la primera mitad es que no.
Las ocho Datasheets se escriben sin campos propios, pero dejan en prosa del registro de procedencia, que ninguna herramienta lee, datos como los volúmenes contratados de Postmark, Resend y Mailgun o el tope del sobrecoste de SendGrid.
En tres casos, el tope de SendGrid, el cupo de destinatarios de MailerSend y la cola de Brevo, PRIME4API devuelve un resultado que el proveedor contradice.

La respuesta a la segunda mitad son las nueve propuestas de la Tabla~\ref{tab:propuestas}: siete amplían lo que el esquema expresa y dos corrigen formas que la guía admite y PRIME4API no lee.
Ocho de ellas salen de un inventario de treinta y cinco elementos no soportados, y la novena, P-05, del fallo del motor ante un plan cuyo precio no es público.
A ellas se suma una rectificación de la guía de ACM2YAML: la guía documenta \texttt{workload} como un mapa, y una Datasheet de ejemplo de su propio repositorio lo escribe como lista (Apartado~\ref{subsec:evolucion:dv01}).

Gavira Sánchez declara que ACM2YAML v0.3 captura dos de las cuatro dimensiones de capacidad del metamodelo, el volumen (las cuotas) y la frecuencia (los \textit{rates}), y deja fuera las peticiones concurrentes y los resultados por petición \citep{gavira2026tfm}.
Los resultados por petición (\textit{Results Per Request}) están documentados en siete APIs del catálogo, y Gavira Sánchez ya recoge su serialización como línea futura \citep{gavira2026tfm}.
Las peticiones concurrentes aparecen en el catálogo en los topes del relé SMTP que reúne F-15: SendGrid admite hasta 10.000 conexiones simultáneas desde un mismo servidor, y MailerSend, cinco en sus planes Free, Hobby y Starter.
Las dos cifras quedan en el registro de procedencia.

Tres propuestas tratan lo que v0.3 sí recoge pero declara con un único valor escalar: la cifra de la cuota y el consumo de cada petición.
Postmark vende seis volúmenes contratados dentro de un mismo plan, donde el esquema admite una sola cifra de cuota (P-01, Apartado~\ref{subsec:evolucion:p01}).

Las otras dos tratan el consumo de cada petición.
Una llamada al lote de Resend admite hasta cien mensajes con hasta cincuenta destinatarios cada uno y consume hasta el producto de las dos cifras, 5.000 correos.
El esquema, en cambio, describe el consumo de una petición con una sola terna de unidad, mínimo y máximo (P-02).
El tope de \texttt{/email/bulk} de Postmark es un tamaño, 50 MB por petición, y el esquema solo admite acotar una petición en su propia unidad, que es el correo (P-03).
Los dos casos se detallan en el Apartado~\ref{subsec:evolucion:p02}.

La capacidad depende además del estado de la suscripción y de la reputación de la cuenta.
Cuando MailerSend pone una cuenta bajo revisión, su envío baja de 120 a 10 peticiones por minuto sin cambiar de plan, y ese eje aparece en seis de las ocho APIs (P-06, Apartado~\ref{subsec:evolucion:p06}).
La guía permite declarar dentro de cada \textit{endpoint} un alias por estado, con la cifra de cada uno, pero obliga a repetirlo en todos los \textit{endpoints}.
Tampoco tiene campo para el disparador que cambia el estado, como la aprobación de la cuenta en MailerSend o su revisión por reputación.

Otras dos propuestas tratan lo que ocurre al superar la cuota.
SendGrid limita el envío mensual a diez veces el volumen contratado, sobrecoste incluido, y el esquema no tiene campo para ese tope (P-07, Apartado~\ref{subsec:evolucion:p07}).
Brevo retiene hasta 1.000 correos por encima de la cuota diaria y deja de entregar los siguientes (P-09, Apartado~\ref{subsec:evolucion:p09}).
F-13, el sobrecoste opcional que no existe en el plan gratuito, se refiere también al exceso sobre la cuota en cinco APIs, y su forma queda esbozada en la Tabla~\ref{tab:sinforma}.

Al margen del exceso sobre la cuota, la cuenta de MailerSend sin aprobar solo puede enviar a dos direcciones distintas, con un cupo que no se repone (P-08, Apartado~\ref{subsec:evolucion:p08}).
P-07, P-08 y P-09 aparecen en una sola API cada una.
Se promueven porque PRIME4API devuelve sobre la Datasheet sin manipular un resultado que el proveedor contradice, y la misma operación responde bien en un caso de control en que el límite le llega escrito, en la Datasheet o en la consulta.

Las dos propuestas restantes son de contrato: el esquema admite el valor que describe al proveedor y PRIME4API falla al leerlo.
En el plan Enterprise de Resend, un límite vacío no dice si no se publica o si no existe, y el motor falla al convertirlo a número (P-04, Apartado~\ref{subsec:evolucion:p04}).
El mismo plan lleva el precio \texttt{custom}, que la guía admite y que impide al motor evaluar el documento entero (P-05, Apartado~\ref{subsec:evolucion:p05}).

Además de las nueve propuestas, la Tabla~\ref{tab:sinforma} recoge doce elementos que aparecen en dos APIs o más y cuya forma está esbozada sin contrastarla con la guía.
Otros cuatro elementos promovidos quedan fuera de las propuestas (Sección~\ref{sec:esquema:fuera}): los resultados por petición, los \textit{add-ons}, los límites que la respuesta no refleja y la superficie de la API que cambia con el plan.
Los \textit{add-ons}, como la IP dedicada, aparecen en las ocho APIs, pero el esquema los deja al canal web y el catálogo solo modela la API REST, una de las interfaces del canal API.

Tampoco hay propuesta para los límites que el proveedor declara y que la respuesta HTTP no refleja.
MailerSend acepta con un \texttt{202} un lote que después rechaza, y Mailchimp Transactional responde con \texttt{200} a una llamada de unos 40 MB cuando declara un tope de 10 MB.
Estos casos plantean qué cuenta como capacidad contractual, y esa cuestión no la resuelve ningún campo nuevo.
La superficie de la API que cambia con el plan ya cabe en v0.3, porque un \textit{endpoint} que un plan no ofrece se omite de ese plan.

Las nueve propuestas y los dieciséis elementos promovidos sin propuesta valen para un dominio y un motor: la recurrencia se cuenta solo en las ocho APIs de correo transaccional, y las Datasheets solo se han evaluado con PRIME4API.
